% Generated draft: requires mathematical and source review.
\documentclass[11pt]{article}
\usepackage[margin=25mm]{geometry}
\usepackage{fontspec}
\setmainfont{Latin Modern Roman}
\usepackage{amsmath,amssymb,mathtools,unicode-math}
\setmathfont{Latin Modern Math}
\usepackage{enumitem,needspace,xurl,hyperref}
\setlength{\parindent}{0pt}
\setlength{\parskip}{5pt}
\setlength{\emergencystretch}{3em}
\newcommand{\N}{\mathbb{N}}
\newcommand{\Z}{\mathbb{Z}}
\newcommand{\Q}{\mathbb{Q}}
\newcommand{\R}{\mathbb{R}}
\newcommand{\C}{\mathbb{C}}
\newcommand{\F}{\mathbb{F}}
\newcommand{\A}{\mathbb{A}}
\newcommand{\PP}{\mathbb{P}}
\newcommand{\E}{\mathbb{E}}
\newcommand{\eps}{\varepsilon}
\newcommand{\dd}{\,\mathrm{d}}
\begin{document}
\section*{UnsolvedMath ID: 2306002}
\textbf{Source problem code:} catalogue entry
\subsection{Definitions and mathematical statement}
Let $S$ be the class of normalized univalent holomorphic functions on $\mathbb D=\{z\in\C:|z|<1\}$, with Taylor expansions
\[
 f(z)=z+\sum_{n=2}^{\infty}a_n(f)z^n.
\]
For $n\geq1$, with $a_1(f)=1$, define the extremal coefficient magnitude
\[
 A_n=\sup_{f\in S}|a_n(f)|.
\]
The source attributes to Hayman the existence of a constant $K_0$ such that
\[
 \lim_{n\to\infty}\frac{A_n}{n}=K_0.
\]
\textbf{Recorded question.} Is $K_0=1$?

The statement quotes the historical estimate $K_0<1.243$, attributed to Milin, as the best known result at the time of that problem text. The supplied later research summary marks the question solved: it attributes $A_n=n$ for every $n$ to de Branges' proof of the Bieberbach conjecture and therefore reports $K_0=1$. This later status is distinguished from the historical ``best known'' wording; no independent proof or literature verification is claimed here.
\subsection{Short English statement}
For normalized univalent functions, let $A_n$ be the largest possible magnitude of the $n$th Taylor coefficient. Hayman's limit question asks whether $A_n/n\to1$. The supplied summary reports this settled by de Branges' theorem, which gives $A_n=n$.
\subsection{Sources}
\begin{itemize}
\item Local UnsolvedMath registry, permanent ID 2306002, source code catalogue entry, ``Research Problems in Function Theory --- Problem 6.2''; supplied status solved, classification SOLVED-IN-LITERATURE, updated 2026-08-06.
\item Hayman, \emph{Research Problems in Function Theory} (2018), Problem 6.2, \url{https://arxiv.org/abs/1809.07200}. The record reports source-TeX extraction and access on 2026-07-29; the 2018 update was retained for status triage. Current status and rights require review; the record describes publicly available arXiv source TeX.
\item The problem text attributes limit existence to Hayman and the historical $1.243$ bound to Milin. The supplied summary attributes the final value to de Branges' proof of the Bieberbach conjecture. No complete original-paper citations are supplied for these attributions.
\end{itemize}
\end{document}
